\documentclass[11pt,a4paper]{article}
\usepackage[utf8]{inputenc}
\usepackage[T1]{fontenc}
\usepackage{xcolor}
\usepackage{graphicx}
\usepackage[colorlinks=true,linkcolor=blue!60!black,citecolor=blue!60!black,urlcolor=blue!60!black]{hyperref}
\usepackage{booktabs}
\usepackage{listings}
\usepackage{longtable}
\usepackage{geometry}
\usepackage{float}
\usepackage{amssymb}
\usepackage[labelfont=bf,font=small,skip=4pt]{caption}
\usepackage{microtype}
\usepackage{parskip}
\usepackage{makecell}
\usepackage{array}
\usepackage{xurl}

\usepackage{tabularx}
\usepackage{booktabs}


\usepackage{lmodern}
\usepackage{xurl}
\usepackage{titlesec}

\renewcommand{\arraystretch}{1.2}
\setlength{\tabcolsep}{6pt}

\setlength{\parskip}{6pt}
\setlength{\parindent}{0pt}

\setlength{\textfloatsep}{10pt}
\setlength{\floatsep}{8pt}
\setlength{\intextsep}{8pt}

\captionsetup{
  justification=raggedright,
  singlelinecheck=false
}

\newcolumntype{Y}{>{\raggedright\arraybackslash}X}

\titlespacing*{\section}{0pt}{14pt}{6pt}
\titlespacing*{\subsection}{0pt}{10pt}{4pt}

\geometry{margin=1in, top=1.05in, bottom=1.05in}
\graphicspath{{docs/figures/report/images/}}

\newcommand{\sectionrule}{\vspace{6pt}\hrule\vspace{10pt}}
\newcommand{\tightitem}{\setlength\itemsep{2pt}\setlength\parskip{0pt}}
\newcommand{\outcomebox}[1]{%
  \begin{center}
  \fbox{\parbox{0.95\linewidth}{#1}}
  \end{center}
}
\newcommand{\demofigure}[4]{%
  \begin{figure}[H]
  \centering
  \IfFileExists{#1}{%
    \includegraphics[width=#2\linewidth]{#1}%
  }{%
    \fbox{\parbox{#2\linewidth}{\centering\vspace{1.3cm}\textit{Placeholder for #3}\\[4pt]\texttt{\detokenize{#1}}\vspace{1.3cm}}}%
  }
  \caption{#4}
  \end{figure}
}

\lstset{
  basicstyle=\ttfamily\small,
  breaklines=true,
  breakatwhitespace=true,
  columns=fullflexible,
  keepspaces=true,
  frame=single,
  framesep=4pt,
  xleftmargin=4pt,
  xrightmargin=4pt
}

\begin{document}

\title{\textbf{Meeting Summarizer}\\[4pt]
\large  Technical Report }
\author{Viswanadha Sai Nissankararao}
\date{March 2026}
\maketitle

\begin{center}
\begin{tabular}{ll}
\textbf{Project:} & Meeting Summarizer (Project 14) \\
\textbf{Dataset:} & SAMSum \\
\textbf{Platform:} & Apple M4 Pro $\cdot$ 24 GB UMA $\cdot$ MPS / BF16 \\
\textbf{Primary champion:} & BART-base + D27: ROUGE-L \textbf{40.12} (target $\geq$ 40) \\
\textbf{Alternate (T5-class):} & FLAN-T5-base E1: ROUGE-L \textbf{42.275} (819-test; \path{google_flan-t5-base_with_speakers_test.json}) \\
\textbf{Alignment doc:} & \path{docs/REPORT_ALIGNMENT.md} (PDF claims vs.\ committed JSON) \\
\end{tabular}
\end{center}

\outcomebox{%
\textbf{Project Outcome.} The project solves abstractive dialogue summarization on the SAMSum benchmark using fine-tuned seq2seq transformers and an interactive Streamlit demo. The primary reported champion is \texttt{facebook/bart-base} trained on speaker-tagged dialogues and decoded with beam search (Config D27: \texttt{num\_beams=5}, \texttt{length\_penalty=1.33}), achieving \textbf{ROUGE-L 40.12} on the full 819-example test set. An alternative path that also exceeds the $\geq 40$ target is \textbf{\texttt{google/flan-t5-base}} full fine-tune (\textbf{42.275} ROUGE-L on the same split). \textbf{\texttt{t5-small}} E1 is \textbf{31.95} ROUGE-L and does \emph{not} meet $\geq 40$ without a different model or checkpoint.
}

\sectionrule

\section{Executive Summary and Success Criteria}

This project builds an end-to-end \textbf{abstractive dialogue summarization system} for multi-speaker conversations. The system takes a chat-style conversation as input and returns a concise summary, along with optional structured outputs in the demo interface.

The evaluation target was explicit: \textbf{ROUGE-L $\geq$ 40.0} on the 819-example SAMSum test set. The criterion is met by \textbf{BART-base} with tuned decoding (\textbf{40.12}) and by \textbf{FLAN-T5-base} full fine-tune (\textbf{42.275}). \textbf{T5-small} E1 reaches \textbf{31.95} and does \emph{not} satisfy the threshold on its own.

Four model families were evaluated: \texttt{facebook/bart-base}, \texttt{t5-small}, \texttt{google/flan-t5-base}, and \texttt{google/pegasus-cnn\_dailymail}. Committed decode-sweep files named \path{t5_decode_*.json} / \path{t5_decoding_sweep_summary.json} refer to the \textbf{FLAN-T5-base} checkpoint unless regenerated for \texttt{t5-small} (see \texttt{hf\_model\_id} in JSON).

\sectionrule

\section{Problem Definition and Scope}

\subsection{Task}

Given a multi-turn, multi-speaker dialogue, generate a concise abstractive summary that preserves the main facts, decisions, and speaker attribution.

\textbf{Example input:}
\begin{lstlisting}
Amanda: I baked cookies. Do you want some?
Jerry: Sure!
Amanda: I'll bring you tomorrow :-)
Jerry: Thanks! Do you know how to make the lemon ones?
Amanda: The biscuits?
Jerry: Yeah.
Amanda: I'll send you the recipe. It's easy!
\end{lstlisting}

\textbf{Example output:}
\begin{lstlisting}
Amanda will bring Jerry cookies tomorrow and send him the recipe for lemon biscuits.
\end{lstlisting}

\subsection{Why This Task Is Difficult}

\begin{itemize}\tightitem
\item Dialogues are informal, compressed, and often omit context.
\item Correct summarization requires \textbf{speaker attribution}, not just topic extraction.
\item The average compression ratio is about \textbf{5.2$\times$}, so omission and hallucination are realistic risks.
\item Pre-trained summarization checkpoints are usually trained on news or web text rather than conversation.
\end{itemize}

\subsection{Success Criterion}

\textbf{Primary target:} ROUGE-L $\geq$ 40.0 on the full SAMSum test split.\\
\textbf{Achieved:} ROUGE-L 40.12 using BART-base + speaker tags + tuned decoding.

\sectionrule

\section{Dataset and Data Quality}

\subsection{Dataset Overview}

\begin{center}
\begin{tabular}{p{0.28\linewidth}p{0.56\linewidth}}
\toprule
\textbf{Property} & \textbf{Value} \\
\midrule
Source & \texttt{knkarthick/samsum} on Hugging Face \\
Reference & Gliwa et al., 2019 \\
Domain & English messenger-style dialogues \\
License & CC BY-NC-ND 4.0 \\
Task type & Abstractive dialogue summarization \\
\bottomrule
\end{tabular}
\end{center}

\begin{center}
\begin{tabular}{lr}
\toprule
Split & Dialogues \\
\midrule
Train & 14,731 \\
Validation & 818 \\
Test & 819 \\
\textbf{Total} & \textbf{16,368} \\
\bottomrule
\end{tabular}
\end{center}

\subsection{Data Quality Evidence}

\begin{itemize}\tightitem
\item \textbf{Leakage check passed:} \texttt{scripts/data\_audit.py} found zero overlap across train, validation, and test IDs.
\item \textbf{Dialogue length:} mean 148.9 tokens, p99 525 tokens.
\item \textbf{Summary length:} mean 28.7 tokens, max 94 tokens.
\item \textbf{Speaker structure:} 73\% of dialogues have exactly two speakers.
\item \textbf{Training implication:} \texttt{max\_source\_length=512} covers nearly all dialogues, so truncation affects only an edge case.
\end{itemize}

\begin{figure}[H]
\centering
\includegraphics[width=0.94\linewidth]{images/dataset_example_dialogue.png}
\caption{Representative SAMSum dialogue example used to illustrate the input structure and target summary style.}
\end{figure}

\subsection{Dataset Risk}

SAMSum is strong for controlled benchmarking, but it is \textbf{constructed rather than transcribed}. That means it is cleaner than real meeting transcripts and may overestimate performance on noisy production conversations.

\sectionrule

\section{Preprocessing and Baselines}

\subsection{Preprocessing Variants}

Three preprocessing variants were evaluated because they directly affect the scoring categories for data handling and experiments:

\begin{itemize}\tightitem
\item \textbf{with\_speakers:} preserve speaker tags such as \texttt{Amanda:}.
\item \textbf{no\_speakers:} strip speaker tags to test the value of explicit attribution.
\item \textbf{split\_speakers:} use sliding-window segmentation for dialogues that exceed 512 tokens.
\end{itemize}

\begin{figure}[H]
\centering
\includegraphics[width=0.94\linewidth]{images/preprocessing_pipeline.png}
\caption{Preprocessing pipeline used to create the three training variants from raw SAMSum dialogues.}
\end{figure}

\subsection{Zero-Shot Baselines}

\begin{center}
\begin{tabular}{lrrrr}
\toprule
Model & ROUGE-1 & ROUGE-2 & ROUGE-L & Test Size \\
\midrule
BART-base & 27.34 & 8.87 & 19.89 & 100 \\
T5-small & 27.60 & 7.63 & 22.19 & 100 \\
PEGASUS & 1.85 & 0.00 & 1.60 & 100 \\
\bottomrule
\end{tabular}
\end{center}

These baselines establish two points: fine-tuning is necessary, and larger models do not automatically transfer well to dialogue summarization.

\sectionrule

\section{Models and Experiment Design}

\subsection{Model Choice}


\begin{center}
\small
\begin{tabularx}{\linewidth}{p{0.28\linewidth}p{0.12\linewidth}X X}
\toprule
\textbf{Model} & \textbf{Params} & \textbf{Role in project} & \textbf{Why included} \\
\midrule
\makecell[l]{\texttt{facebook/bart-base}} & 139M & Primary model & Strong encoder-decoder baseline for abstractive summarization \\

\makecell[l]{\texttt{t5-small}} & 60M & Compact comparison baseline & Tests a smaller text-to-text model under the same task \\

\makecell[l]{\texttt{google/pegasus-}\\\texttt{cnn\_dailymail}} & 568M & Transfer stress test & Tests cross-domain transfer from news summarization \\

\bottomrule
\end{tabularx}
\end{center}



\subsection{Experiment Matrix}

\begin{center}
\small
\begin{longtable}{p{0.08\linewidth}p{0.29\linewidth}p{0.12\linewidth}p{0.13\linewidth}p{0.28\linewidth}}
\toprule
\textbf{ID} & \textbf{Change tested} & \textbf{ROUGE-L} & \textbf{Cost} & \textbf{Conclusion} \\
\midrule
\endfirsthead
\toprule
\textbf{ID} & \textbf{Change tested} & \textbf{ROUGE-L} & \textbf{Cost} & \textbf{Conclusion} \\
\midrule
\endhead
E0 & Zero-shot baselines & 19.89--22.19 & eval only & Fine-tuning is required \\
E1 & BART vs T5 full fine-tuning & 39.85 vs 31.95 & 72.4 vs 35.1 min & BART is clearly stronger \\
E2 & Remove speaker tags & 33.23 & standard train & Speaker attribution is crucial \\
E3 & Tune decoding hyperparameters & \textbf{40.12} & no retraining & Best final result \\
E4 & Faithfulness evaluation & n/a & eval only & ROUGE alone is insufficient \\
E5 & LoRA instead of full fine-tune & 37.59 & 54.7 min & Efficient but lower quality \\
E6 & Sliding-window handling & 39.08 & standard train & Extra complexity not justified \\
E7 & PEGASUS transfer & 1.15 & 1 epoch & Severe domain mismatch \\
E8 & Longer training schedule & 38.46 & 259.6 min & Slower and worse than default \\
\bottomrule
\end{longtable}
\end{center}

\sectionrule

\section{Training Setup and Efficiency}

\subsection{Core Training Configuration}

\begin{center}
\begin{tabular}{p{0.30\linewidth}p{0.22\linewidth}p{0.33\linewidth}}
\toprule
\textbf{Setting} & \textbf{Value} & \textbf{Why it matters} \\
\midrule
Max source length & 512 & Covers nearly all dialogues \\
Max target length & 128 & Above the observed max summary length \\
Batch size & 8 & Safe on 24 GB UMA with BF16 \\
Learning rate & 5e-5 & Standard fine-tuning configuration \\
Epochs & 5 & Best schedule among tested defaults \\
Early stopping & patience = 2 & Reduces wasted computation \\
Default decoding & beam = 4, lp = 1.0 & Training-time baseline \\
Best decoding & beam = 5, lp = 1.33 & Final quality gain without retraining \\
\bottomrule
\end{tabular}
\end{center}

\subsection{Efficiency Summary}

\begin{center}
\small
\begin{tabular}{p{0.28\linewidth}p{0.16\linewidth}p{0.16\linewidth}p{0.26\linewidth}}
\toprule
\textbf{Configuration} & \textbf{Train time} & \textbf{Key metric} & \textbf{Efficiency takeaway} \\
\midrule
T5-small (E1) & 35.1 min & ROUGE-L 31.95 & Faster but well below target \\
BART-base (E1) & 72.4 min & ROUGE-L 39.85 & Best full fine-tune baseline \\
BART-base + D27 (E3) & retrain not needed & ROUGE-L 40.12 & Best overall result \\
BART-base LoRA (E5) & 54.7 min & ROUGE-L 37.59 & Useful speed/quality compromise \\
BART extended (E8) & 259.6 min & ROUGE-L 38.46 & Much slower and worse \\
\bottomrule
\end{tabular}
\end{center}

Training logs show post-training MPS memory near \textbf{16.96 GB}, which is consistent with the selected batch size and BF16 configuration on Apple Silicon.

\sectionrule

\section{Evaluation Methodology}

\subsection{Primary Metrics}

\begin{itemize}\tightitem
\item \textbf{ROUGE-1:} unigram overlap.
\item \textbf{ROUGE-2:} bigram overlap.
\item \textbf{ROUGE-L:} longest common subsequence; this is the \textbf{primary selection metric}.
\end{itemize}

All ROUGE values are reported as macro-averaged F-measures $\times 100$ over the full 819-example test set unless otherwise noted.

\subsection{Supporting Evaluation}

The report uses three supporting checks beyond ROUGE:

\begin{itemize}\tightitem
\item \textbf{Bootstrap confidence intervals} for the main architecture comparison.
\item \textbf{Faithfulness metrics} including entity hallucination, speaker preservation, and NLI entailment.
\item \textbf{Manual error analysis} on 20 examples for qualitative failure inspection.
\end{itemize}

\textbf{Human evaluation scope.} This submission includes a completed \textbf{Task 3 human-eval pilot}: \path{results/steering/facebook_bart-base_with_speakers_human_eval_rated_backup.csv} with protocol in \path{results/steering/human_eval_rubric.md}, summarized in \path{results/metrics/task3_human_eval_pilot_summary.json} (200 rated rows = 50 samples $\times$ 4 $\alpha$ settings). Task 4 coherence remains manual CSV-based unless those sheets are also filled and summarized. See \path{docs/REPORT_ALIGNMENT.md}.

\begin{center}
\begin{tabular}{p{0.30\linewidth}p{0.22\linewidth}p{0.34\linewidth}}
\toprule
\textbf{Metric} & \textbf{Value} & \textbf{Interpretation} \\
\midrule
Entity hallucination rate & 10.1\% (83/819) & Roughly 1 in 10 outputs invents an entity \\
Speaker preservation & 75.5\% & Most summaries preserve speaker identity \\
NLI faithfulness & 0.308 & Moderate entailment; useful but insufficient alone \\
\bottomrule
\end{tabular}
\end{center}

\sectionrule

\section{Results and Best Model Selection}

\subsection{Main Results}

\begin{figure}[H]
\centering
\includegraphics[width=0.94\linewidth]{images/model_performance_comparison.png}
\caption{Zero-shot and fine-tuned model comparison with uncertainty on the fine-tuned systems.}
\end{figure}

\begin{figure}[H]
\centering
\includegraphics[width=0.94\linewidth]{images/experiment_rougel_summary.png}
\caption{ROUGE-L across the reported experiments. The final system crosses the project target.}
\end{figure}

\subsection{Key Findings}

\begin{itemize}\tightitem
\item \textbf{Architecture:} BART-base outperformed T5-small by \textbf{+7.90 ROUGE-L} after fine-tuning.
\item \textbf{Speaker tags:} removing them reduced ROUGE-L from 39.85 to \textbf{33.23}, making this the most important preprocessing decision.
\item \textbf{Decoding:} D27 improved ROUGE-L from 39.92 to \textbf{40.12} without retraining.
\item \textbf{Efficiency:} LoRA retained useful quality (37.59) with less cost than full fine-tuning.
\item \textbf{Longer training:} E8 was slower and worse, so extra epochs were not justified.
\end{itemize}

\subsection{Best Model Selection}

\begin{center}
\begin{tabular}{p{0.28\linewidth}p{0.54\linewidth}}
\toprule
\textbf{Property} & \textbf{Best configuration} \\
\midrule
Model & \texttt{facebook/bart-base} \\
Preprocessing & \texttt{with\_speakers} \\
Decoding & \texttt{num\_beams=5}, \texttt{length\_penalty=1.33} (D27) \\
ROUGE-1 / ROUGE-2 / ROUGE-L & 48.48 / 23.55 / \textbf{40.12} \\
Inference latency & 200.6 ms/sample \\
Reason selected & Highest test ROUGE-L while remaining reproducible on local hardware \\
\bottomrule
\end{tabular}
\end{center}

\subsection{FLAN-T5-base (alternative $\geq 40$ path)}

Full fine-tuning of \texttt{google/flan-t5-base} with speaker-tagged preprocessing yields \textbf{ROUGE-L 42.275} on the 819-example test set (\path{results/metrics/google_flan-t5-base_with_speakers_test.json}), exceeding the project target without relying on the BART decoding sweep. This does \emph{not} change the fact that \texttt{t5-small} E1 remains near \textbf{31.95}; always name the checkpoint when citing a score.

\sectionrule

\section{Error Analysis and Key Insights}

\subsection{Error Distribution}

\begin{figure}[H]
\centering
\includegraphics[width=0.92\linewidth]{images/error_distribution.png}
\caption{Manual error analysis for a 20-example sample from the champion model outputs.}
\end{figure}

\begin{center}
\begin{tabular}{lrr}
\toprule
Category & Count & Share \\
\midrule
Correct & 4 & 20\% \\
Partial & 8 & 40\% \\
Hallucination & 6 & 30\% \\
Over-generic & 2 & 10\% \\
Truncated & 0 & 0\% \\
\bottomrule
\end{tabular}
\end{center}

\subsection{Representative Failure Cases}

\begin{center}
\small
\begin{tabular}{p{0.22\linewidth}p{0.28\linewidth}p{0.36\linewidth}}
\toprule
\textbf{Failure type} & \textbf{Observed output} & \textbf{Why it matters} \\
\midrule
Hallucination & ``Pogba scored the first goal of the season'' & The dialogue discussed form, not a goal; this is fabricated content \\
Action inversion & ``Randolph will buy earplugs for Maya'' & The vocabulary is correct, but actor and recipient are reversed \\
Missing action detail & ``Tom and Ben will meet at 2 pm'' & Logistics are preserved, but the instruction to bring papers is lost \\
\bottomrule
\end{tabular}
\end{center}

\subsection{Three High-Value Insights}

\begin{itemize}\tightitem
\item \textbf{Speaker tags are the most impactful preprocessing choice.}
\item \textbf{Decoding optimisation is free performance gain} once the model is trained.
\item \textbf{Faithfulness remains the main weakness} even when ROUGE-L is strong.
\end{itemize}

\sectionrule

\section{Project 3: Advanced Capabilities (Tasks 1--5)}

The repository extends the core E0--E8 experiments with five research tracks. The following paragraphs are aligned with \path{docs/REPORT_ALIGNMENT.md} and committed artifacts under \path{results/metrics/}; they take precedence over informal slide wording.

\subsection{Task 1: Attention analysis}

\path{scripts/task1_attention_patterns.py} and related utilities provide attention visualizations and diagnostics. Interpretations should respect the script's distinction between teacher-forced attention and production beam decoding.

\subsection{Task 2: Quantization and latency benchmarks}

\textbf{Quantization.} Models are exported for \textbf{CTranslate2}; folder names such as \texttt{Q4\_K\_M} / \texttt{Q5\_K\_M} / \texttt{Q8\_0} are nominal mappings in the project layout --- \emph{not} a claim of llama.cpp GGUF conversion. Evidence: \path{task2_quantization_manifest.json}, \path{task2_benchmark_table.json}.

\textbf{Streaming vs.\ batch.} \path{scripts/task2_benchmark.py} implements a rolling incremental summarizer for ablation. Latencies in \path{task2_streaming_vs_batch.json} compare two \emph{harness} modes; streaming can report much higher end-to-end time than batch. Treat these numbers as \textbf{diagnostic}, not as a production ``real-time'' SLA or sub-second live meeting latency; batch inference remains the default recommendation for wall-clock work.

\subsection{Task 3: Steering and human evaluation}

Activation steering experiments use automated metrics plus a completed \textbf{human-eval pilot} (Task 3) with protocol-backed ratings and summary JSON under \path{results/metrics/}.

\subsection{Task 4: Adversarial robustness}

Adversarial training and LoRA retraining are evaluated against \path{task4_robustness_comparison.json}. The aggregate micro headline \texttt{robustness\_gain} is approximately \textbf{$-0.07$} --- i.e.\ \textbf{no} aggregate ROUGE improvement on the headline metric. Per-pattern values in \texttt{robustness\_gain\_by\_pattern} show trade-offs (some patterns up, others flat or down). Do not describe the retrain as uniformly improving adversarial robustness without citing these fields.

\subsection{Task 5: Structured JSON output}

Structured evaluation in \path{task5_structured_output.json} (see field \texttt{n\_samples}, typically 819 at full test) distinguishes: \textbf{(i)} \texttt{strict\_generative\_json\_rate} --- success of \texttt{json.loads} on the raw model string (after whitespace repair only); \textbf{(ii)} \texttt{salvaged\_json\_rate} --- deterministic salvage of JSON-ish T5 text; \textbf{(iii)} \texttt{generative\_native\_json\_rate} = strict + salvage. Committed runs often show \textbf{strict = 0} with \textbf{salvage-mediated} recovery; reports should not imply ``the model always outputs valid JSON'' without naming the column. A valid API envelope via post-processing (\texttt{structured\_dict\_from\_model\_output}) is \emph{not} the same as strict generative JSON.

\sectionrule

\section{Demo Application}

\subsection{What the Demo Shows}

The Streamlit demo provides explicit evidence that the project is not only a training pipeline but also a usable application. The interface supports:

\begin{itemize}\tightitem
\item dialogue input through a text area,
\item model selection from available checkpoints,
\item decoding controls such as beam width and length penalty,
\item generated summary output,
\item extracted action items,
\item named entities,
\item generation metadata including model name, device, and latency.
\end{itemize}

\subsection{ Screenshots  }

\demofigure{images/streamlit_ui_input.png}{0.96}{UI input screenshot}{Demo screenshot before generation: the dialogue input area is visible together with enough surrounding UI context to show how a user starts the workflow.}

\demofigure{images/streamlit_full_results.png}{0.96}{UI output screenshot}{Demo screenshot after generation: the summary output is visible together with at least one structured output such as action items, named entities, or the latency/model information block.}

\subsection{Qualitative Demo Examples}

\begin{center}
\small
\begin{tabular}{p{0.18\linewidth}p{0.31\linewidth}p{0.25\linewidth}p{0.18\linewidth}}
\toprule
\textbf{Example} & \textbf{Input type} & \textbf{Expected visible output} & \textbf{Why it helps scoring} \\
\midrule
A & Short coordination dialogue & Faithful summary + correct speaker action & Shows baseline usability \\
B & Longer planning dialogue & Summary + action items + entities & Shows richer structured output \\
\bottomrule
\end{tabular}
\end{center}

\sectionrule

\section{Reproducibility}

\subsection{Minimal Reproduction Path}

\begin{lstlisting}
# 1) Environment
python3 -m venv .venv
source .venv/bin/activate
pip install -r requirements.txt

# 2) Data and preprocessing
python3 scripts/predownload_assets.py
python3 scripts/data_audit.py
python3 scripts/preprocess.py --variants with_speakers no_speakers split_speakers

# 3) Training and evaluation
python3 scripts/train.py
python3 scripts/decoding_ablation.py
python3 scripts/evaluate_faithfulness.py
python3 scripts/bootstrap_ci.py

# 4) Demo
bash scripts/run_app.sh
\end{lstlisting}

\subsection{Critical Files}


\begin{center}
\small
\begin{tabular}{p{0.32\linewidth}p{0.58\linewidth}}
\toprule
\textbf{Path} & \textbf{Purpose} \\
\midrule
\path{config.yaml} & primary training and decoding configuration \\
\path{config_extended.yaml} & extended-schedule experiment (E8) \\
\path{scripts/train.py} & main fine-tuning pipeline \\
\path{scripts/decoding_ablation.py} & decoding sweep used for D27 selection \\
\path{results/metrics/} & all reported quantitative outputs \\
\path{docs/REPORT_ALIGNMENT.md} & checklist: JSON definitions, forbidden claims, Task 2/4/5 wording \\
\path{models/best/} & best saved checkpoints \\
\path{scripts/app.py} & Streamlit demo application \\
\bottomrule
\end{tabular}
\end{center}


All experiments use \texttt{seed=42}, and the report only cites metrics that are already committed in repository outputs.

\textbf{PDF alignment.} After editing metrics or prose, regenerate this PDF (or update the \texttt{.tex} source) so claims match \path{docs/REPORT_ALIGNMENT.md} and the JSON files listed there.

\sectionrule

\section{Conclusion}

\begin{itemize}\tightitem
\item The project \textbf{met its success criterion}, achieving ROUGE-L 40.12 (BART + D27) and also \textbf{42.275} with FLAN-T5-base E1 on the full SAMSum test split.
\item The primary champion system is \textbf{BART-base with speaker tags and tuned beam-search decoding}; FLAN-T5-base provides an alternative path above the same threshold.
\item The strongest empirical insight is that \textbf{speaker attribution is essential} for dialogue summarization quality.
\item The main remaining weakness is \textbf{faithfulness}, especially hallucination and action-direction inversion.
\item \textbf{Project 3 limitations:} Task~2 streaming figures are diagnostic, not production SLAs; Task~4 aggregate robustness gain is not positive on the headline metric; Task~5 structured output is salvage-aware --- not raw strict JSON from the generator alone; Task~3 human evaluation is pilot-complete, while Task~4 coherence sheets remain manual unless additionally filled.
\item The project is \textbf{reproducible and demo-ready} on consumer Apple Silicon hardware.
\end{itemize}

\vspace{6pt}


\end{document}
